\section{Features}

\subsection{LaTeX editor}

EcTeX provides a specialized editor for \texttt{.tex} files
with syntax highlighting. It also offers productivity
shortcuts for frequently used LaTeX commands.

For example, the following commands are available by
default:

\begin{itemize}[noitemsep]
    \item \textbf{Ctrl+B}: Insert \verb|\textbf{}|.
    \item \textbf{Ctrl+I}: Insert \verb|\textit{}|.
    \item \textbf{Ctrl+U}: Insert \verb|\underline{}|.
    \item \textbf{Ctrl+Shift+E}: Insert \verb|\emph{}|.
    \item \textbf{Ctrl+Shift+M}: Insert inline mathematics.
\end{itemize}

When text is selected, formatting commands wrap the
selection. Otherwise, an empty command pair is inserted
and the caret is placed inside it.

All shortcuts can be customized through the standard
Eclipse key-binding preferences.

\subsection{Integrated PDF preview}

The PDF Preview view displays the compiled document
next to the source editor. It automatically follows the
active LaTeX file and loads the corresponding PDF.

The preview supports:
\begin{itemize}[noitemsep]
    \item Mouse-wheel scrolling.
    \item Zooming using Ctrl+mouse wheel or keyboard
    shortcuts.
    \item Fit-to-width and zoom commands.
    \item Highlighting of synchronized source locations.
\end{itemize}

PDF rendering is provided through Apache PDFBox
\cite{pdfbox}.

\subsection{Build integration}

EcTeX integrates the standard command-line tools of a TeX
distribution: a LaTeX engine (\texttt{pdflatex} by default, or
\texttt{xelatex} or \texttt{lualatex}) and a bibliography tool
(\texttt{bibtex} by default, or \texttt{biber}). Both are chosen
in the preferences.

The default build commands are:

\begin{center}
\begin{tabular}{ll}
\toprule
\textbf{Shortcut} & \textbf{Operation} \\
\midrule
Ctrl+F2 & Build PDF (one LaTeX pass) \\
Ctrl+F3 & Run the bibliography tool \\
Ctrl+F5 & Build All \\
\bottomrule
\end{tabular}
\end{center}

A LaTeX pass executes:

\begin{lstlisting}
pdflatex -interaction=nonstopmode \
         -file-line-error \
         -synctex=1
\end{lstlisting}

Compiler output is displayed in the EcTeX Build console.
LaTeX errors are also reported through Eclipse problem
markers.

For documents containing bibliographic references,
the recommended compilation sequence is:

\begin{enumerate}[noitemsep]
    \item Run LaTeX.
    \item Run the bibliography tool.
    \item Run LaTeX again.
    \item Run LaTeX once more.
\end{enumerate}

Build All performs exactly this sequence and stops at the
first error. For documents without a bibliography, the
preferences offer the shorter sequence of two LaTeX passes.

\subsection{Documents with several files}

Larger documents are often split over several files, as this document is.
In a LaTeX project (a project converted with
\texttt{Configure $\rightarrow$ Convert to LaTeX Project}, or created with
the New LaTeX Project wizard), one file can be marked as the \emph{main file}
by right-clicking it in the Project Explorer and choosing
\texttt{EcTeX $\rightarrow$ Set as Main File}. The setting is stored per
project, the file is shown with the suffix \texttt{[main file]}, and the
setting is removed again with \texttt{Clear Main File} on the main file or
on the project.

With a main file set, working in any file of the project behaves as if the
main file were open:

\begin{itemize}[noitemsep]
    \item Build and BibTeX run on the main file, after saving the
    project's unsaved files.
    \item The PDF preview shows the PDF of the main file.
    \item Forward synchronization works from every file, and inverse
    synchronization opens the file that contains the clicked text.
    \item Compilation errors are reported on the file that contains them.
\end{itemize}

To try it, set \texttt{ectex-demo-multifile.tex} as the main file and build
from any file in the \texttt{sections} folder.
